\documentclass[a4paper,10pt]{article}
\usepackage[margin=0.5in]{geometry}
\usepackage[hidelinks]{hyperref}
\usepackage{enumitem}
\usepackage{titlesec}
\usepackage{parskip}
\usepackage{array}
\usepackage{tabularx}
\usepackage{multicol}
\usepackage{booktabs}

% ---------- Typography & Spacing Tweaks ----------
\renewcommand{\baselinestretch}{1.03}
\setlist[itemize]{leftmargin=1.5em, itemsep=0pt, parsep=0pt, topsep=1pt}

% Section style: compact spacing and a subtle rule
\titleformat{\section}{\large\bfseries}{}{0em}{}[\vspace{2pt}\hrule height 0.4pt \vspace{6pt}]
\titlespacing*{\section}{0pt}{8pt}{6pt}

% Small helper to tighten vertical gaps
\newcommand{\tightvspace}{\vspace{2pt}}

\begin{document}

% ---------- Header ----------
\begin{center}
    {\LARGE \textbf{Prabesh Kunwar}} \\
    \tightvspace
    \href{mailto:prabeshkunwar12@gmail.com}{prabeshkunwar12@gmail.com} \textbar\ +1 902 598 9453 \textbar\ 
    \href{https://www.linkedin.com/in/prabesh-kunwar/}{LinkedIn} \textbar\ 
    \href{https://github.com/prabeshkunwar12}{GitHub} \textbar\ 
    \href{https://prabeshkunwar12.github.io/portfolio/}{Portfolio}
\end{center}

% ---------- Summary ----------
\noindent \textbf{Prospective MSc (Thesis)} student in Computer Science with a strong background in \textbf{software engineering} and a growing focus on \textbf{quantum computing research}. While I have not completed a formal undergraduate thesis, I have independently pursued research-oriented work in \textbf{quantum algorithms} and \textbf{quantum programming} using \textbf{IBM Qiskit}. This includes updating legacy quantum codebases to current APIs, implementing core quantum algorithms, executing circuits on \textbf{real IBM Quantum hardware}, and analyzing algorithm performance under noisy conditions typical of \textbf{NISQ devices}.

I hold a \textbf{Bachelor of Science in Computer Science} and was named to the \textbf{Dean’s Honours List (2021–2022)} for academic excellence. Alongside my academic preparation, I bring substantial technical experience as a \textbf{software engineer and technical lead}, working on large-scale, real-time systems that integrate software, hardware, and networking. This experience has strengthened my skills in \textbf{systems design}, \textbf{debugging complex systems}, and \textbf{rigorous problem solving skills} directly transferable to research.

I am applying to the thesis-based MSc program to gain \textbf{formal research training}, work under \textbf{faculty supervision}, and develop publishable research in quantum computing and related areas.

\section*{Research Interests}
Quantum Computing; Quantum Algorithms; Noisy Intermediate-Scale Quantum (NISQ) Systems; 
Quantum Programming Frameworks; Computational Modeling; Machine Learning

% ---------- Experience ----------
\section*{Professional Experience}

\textbf{Aerosports Trampoline Park} \textbar\ St. Catherines, Ontario \\
\textit{Technical Lead} \hfill Oct 2024 -- Present \\
\textit{Game Developer / Full-Stack Developer} \hfill Jun 2024 -- Oct 2024
\begin{itemize}
    \item Contributed to the design, implementation, and deployment of \textbf{interactive real-time gaming systems} combining software, hardware controllers, and low-latency network communication.
    
    \item Implemented interactive floor-tile and target-based games using \textbf{.NET} and \textbf{UDP-based sensor communication}, focusing on real-time input handling, synchronization, and system reliability.
    
    \item Developed and maintained backend services using \textbf{Next.js} and \textbf{Express.js}, integrated with WAN-connected databases and administrative dashboards for game, player, and system management.
    
    \item Designed and maintained \textbf{centralized logging, debugging, and simulation tools} to diagnose hardware and software faults, enabling systematic testing and reproducibility across deployments.
    
    \item As Technical Lead, directed the end-to-end technical architecture of the PixelGames platform, translating operational requirements into \textbf{scalable and generalizable system designs} deployable across multiple physical locations.
    
    \item Led the transition from local-only installations to \textbf{region-based, internet-accessible deployments}, enabling centralized configuration, monitoring, and remote updates.
    
    \item Coordinated with construction, electrical, and installation teams to standardize hardware, networking, and software configurations, ensuring consistent system behavior across sites.
    
    \item Oversaw multi-site deployments and long-term system maintenance, strengthening expertise in \textbf{distributed systems, real-time computation, debugging complex systems}, and technical leadership.
\end{itemize}
% ---------- Projects ----------

\section*{Publications \& Thesis}
None to date.

% ---------- Research & Academic Projects ----------
\section*{Research \& Academic Projects}

\textbf{Independent Quantum Computing Project} \\
\textit{The Complete Quantum Computing Course for Beginners} \hfill 2023 -- 2024 \\
\href{https://github.com/prabeshkunwar12/The-Complete-Quantum-Computing-Course-for-Beginners-}{GitHub Repository}
\begin{itemize}
    \item Conducted an in-depth, hands-on study of \textbf{quantum computing fundamentals and quantum algorithms} as part of a structured certification program.
    \item Updated and refactored legacy quantum computing codebases to align with \textbf{current IBM Qiskit APIs}, ensuring correctness and reproducibility under evolving frameworks.
    \item Implemented and analyzed core quantum algorithms, including state preparation, measurement, and algorithmic primitives, using both \textbf{simulators and real IBM Quantum hardware}.
    \item Examined algorithm behavior under \textbf{noisy intermediate-scale quantum (NISQ)} conditions, documenting hardware-specific errors, noise effects, and execution limitations.
    \item Produced detailed explanations and documentation accompanying code implementations, reinforcing conceptual understanding and research-oriented reasoning.
\end{itemize}

\textbf{Life Expectancy Prediction Using Machine Learning} \\
\textit{Data Science and Statistical Modeling Project} \hfill 2022 \\
\href{https://github.com/prabeshkunwar12/lifeExpectancy}{GitHub Repository}
\begin{itemize}
    \item Conducted a cross-country empirical study analyzing the relationship between \textbf{socioeconomic, health, and demographic factors} and life expectancy using WHO datasets.
    \item Performed data cleaning, validation, and exploratory analysis to address missing values, outliers, and inconsistencies in real-world data.
    \item Applied and evaluated \textbf{Polynomial Regression} and \textbf{Random Forest Regressor} models, using cross-validation and hyperparameter tuning to compare predictive performance.
    \item Analyzed model accuracy using \textbf{RMSE and $R^2$ metrics}, identifying key predictors such as schooling, adult mortality, and income composition.
    \item Documented methodology and findings in a formal technical report following an \textbf{IEEE-style academic format}.
\end{itemize}

\textbf{Hotel Rating Microservice Application} \\
\textit{Distributed Systems Project} \hfill 2023 \\
\href{https://github.com/prabeshkunwar12/micro_yt}{GitHub Repository}
\begin{itemize}
    \item Designed and implemented a \textbf{microservices-based architecture} using Spring Boot, emphasizing modularity, scalability, and service discovery.
    \item Integrated \textbf{Netflix Eureka}, \textbf{OpenFeign}, and API Gateway patterns to support inter-service communication and fault tolerance.
    \item Implemented centralized configuration management, authentication, and database persistence using relational and NoSQL databases.
    \item Applied CI/CD workflows to support automated testing, deployment, and reproducibility of distributed systems.
\end{itemize}


% ---------- Education & Achievements (Two Columns) ----------
\section*{Education \& Achievements}

\noindent
\begin{tabularx}{\textwidth}{@{}>{\raggedright\arraybackslash}X >{\raggedright\arraybackslash}X@{}}
\textbf{Education} & \textbf{Achievements \& Certifications} \\
\addlinespace[4pt]

\textbf{University of Prince Edward Island} \par
Charlottetown, PE \par
Bachelor of Science, Major in Computer Science \par
Teaching: Grader -- CS 3420 (Computer Communications), CS 2620 (Programming Languages); 
Lab Assistant -- CS 1910 (Computer Science I)
&
\begin{itemize}[leftmargin=1.2em, itemsep=2pt, topsep=2pt]
  \item \textbf{Dean’s Honours List (2021--2022)}, University of Prince Edward Island

  \item \href{https://coursera.org/share/5969e8ed899773d422e5701ac559e1c1}
  {\textbf{The Complete Quantum Computing Course for Beginners}} — Coursera

  \item \href{https://coursera.org/share/158ce8078de18f7b008995f73c2661be}
  {\textbf{Quantum Computing with Qiskit and Advanced Algorithms}} — Coursera

  \item \href{https://www.coursera.org/account/accomplishments/records/JRF854Y0WN9L}
  {\textbf{Python Programming for Quantum Computing}} — Coursera

  \item \href{https://www.coursera.org/account/accomplishments/verify/UUQLY52S2333}
  {\textbf{Mathematical Foundations and Quantum Mechanics Essentials}} — Coursera

  \item \href{https://www.coursera.org/account/accomplishments/verify/A3JHFH1FVXWV}
  {\textbf{Foundations of Project Management}} — Google, Coursera
\end{itemize}
\end{tabularx}



\end{document}
